\section{Peta Konsep Logika Matematika}

Sebagai ringkasan visual dari konteks kurikulum OBE untuk mata kuliah ini, peta konsep berikut berfungsi sebagai orientasi bagi mahasiswa untuk memahami alur keseluruhan materi Logika Matematika dalam buku ini. Alur pembelajaran mengikuti struktur standar yang digunakan oleh Open Logic Project dan Discrete Mathematics Levin: dimulai dari logika proposisional (proposisi, operator, tabel kebenaran), dilanjutkan ke ekuivalensi dan tautologi, kemudian logika predikat dengan kuantifier, metode inferensi dan pembuktian, induksi matematika, dan terakhir pemodelan logika predikat orde pertama \cite{openlogic,levin_discrete}.

Bab III dan IV membahas logika proposisional secara mendalam; Bab V memperluas ke logika predikat; Bab VI dan VII fokus pada aturan inferensi dan metode bukti termasuk induksi; Bab VIII memperkenalkan pemodelan logika predikat orde pertama sebagai aplikasi praktis. Memahami keterkaitan antarbab membantu mahasiswa menempatkan setiap topik dalam konteks keseluruhan dan melihat bagaimana konsep-konsep tersebut saling membangun.

Berikut peta konsep topik utama per bab:

\begin{center}
\begin{tabular}{|c|l|}
\hline
\textbf{Bab} & \textbf{Topik Utama} \\
\hline
II & Landasan Teori: Logika dan Informatika \\
III & Proposisi $\rightarrow$ Operator $\rightarrow$ Tabel Kebenaran \\
IV & Ekuivalensi $\rightarrow$ Tautologi $\rightarrow$ Translasi \\
V & Predikat $\rightarrow$ Kuantifier $\rightarrow$ Negasi \\
VI & Inferensi $\rightarrow$ Bukti Langsung $\rightarrow$ Kontraposisi \\
VII & Induksi Matematika \\
VIII & Pemodelan FOL: Fakta, Aturan, Inferensi \\
\hline
\end{tabular}
\end{center}

Gunakan peta konsep ini sebagai acuan saat mempelajari setiap bab; rujuk kembali tabel di atas untuk mengingat posisi topik yang sedang dipelajari dalam keseluruhan alur. Hal ini membantu menghindari kehilangan konteks dan memperkuat pemahaman koneksi antar konsep.
